\documentclass[11pt, a4paper]{article}

\usepackage[margin=2.5cm]{geometry}
\usepackage{booktabs}
\usepackage{tabularx}
\usepackage{multirow}
\usepackage{amsmath}
\usepackage{graphicx}
\usepackage{hyperref}
\usepackage{xcolor}
\usepackage{caption}
\usepackage{subcaption}
\usepackage{float}
\usepackage{parskip}
\usepackage{microtype}
\usepackage{fancyhdr}
\usepackage{titlesec}
\usepackage[T1]{fontenc}
\usepackage{lmodern}
\usepackage{amssymb}

\hypersetup{
    colorlinks=true,
    linkcolor=black,
    urlcolor=blue,
    citecolor=black
}

\pagestyle{fancy}
\fancyhf{}
\rhead{Gas Injury Detection — Model Study}
\lhead{\today}
\rfoot{Page \thepage}

\titleformat{\section}{\large\bfseries}{}{0em}{}[\titlerule]
\titleformat{\subsection}{\normalsize\bfseries}{}{0em}{}

% ── Helpers ───────────────────────────────────────────────────────────────────
\newcommand{\acc}[1]{\textbf{#1}}
\newcommand{\pass}{\textcolor{green!60!black}{\textbf{PASS}}}
\newcommand{\fail}{\textcolor{red}{\textbf{FAIL}}}

% ─────────────────────────────────────────────────────────────────────────────
\begin{document}

\begin{titlepage}
    \centering
    \vspace*{2cm}
    {\LARGE\bfseries Gas Injury Detection via Multi-Channel\\[0.4em]
    Gas Sensor Array and Machine Learning\par}
    \vspace{1.5cm}
    {\large Model Study Report\par}
    \vspace{1cm}
    {\normalsize Brent Hoang\par}
    \vspace{0.5cm}
    {\normalsize \today\par}
    \vfill
\end{titlepage}

% =============================================================================
\section{Dataset}
% =============================================================================

Data was collected across 18 sessions using an ESP32-based sensor array with
6 electrochemical channels: VOC, NH$_3$, HCHO, H$_2$S, EtOH, and RH.
16 sessions were used for training and 2 were held out as unseen test data.

Each session was manually labelled into three classes:
\textbf{Baseline} (no sample), \textbf{Sweat}, and \textbf{Blood}.
Blood sessions contain only baseline and blood regions; sweat sessions contain
only baseline and sweat regions.

\begin{table}[H]
\centering
\caption{Dataset summary.}
\begin{tabular}{lr}
\toprule
\textbf{Property} & \textbf{Count} \\
\midrule
Total sessions              & 18 (10 blood, 8 sweat) \\
Training sessions           & 16 \\
Held-out test sessions      & 2 (1 blood, 1 sweat) \\
Training windows (balanced) & 572 \\
Test windows                & 125 \\
\bottomrule
\end{tabular}
\end{table}

% =============================================================================
\section{Methodology}
% =============================================================================

\subsection{Feature Engineering}

Causal helper columns were computed from the 6 raw channels to give models
access to temporal context without look-ahead bias:

\begin{table}[H]
\centering
\caption{Engineered feature groups (per channel, 6 channels total).}
\begin{tabular}{llr}
\toprule
\textbf{Group} & \textbf{Description} & \textbf{Count} \\
\midrule
Raw signal                    & Direct sensor reading              &  6 \\
Rate of change (\texttt{\_roc})   & \texttt{diff()} — instantaneous trend  &  6 \\
Acceleration (\texttt{\_acc})     & \texttt{diff()} of roc             &  6 \\
Rolling mean (15 / 30 / 60 s) & Smoothed level at 3 timescales     & 18 \\
Rolling std  (15 / 30 / 60 s) & Local volatility                   & 18 \\
Rolling roc  (15 / 30 / 60 s) & Smoothed absolute rate of change   & 18 \\
\midrule
\textbf{Total} & & \textbf{72} \\
\bottomrule
\end{tabular}
\end{table}

\textbf{Random Forest / XGBoost features:} each 60\,s non-overlapping window is
summarised into 4 aggregates (mean, std, max, slope) $\times$ 72 columns =
\textbf{288 features} per window.

\textbf{CNN / TCN features:} a 30\,s sliding window (stride 5\,s) retains the
raw time-series with \textbf{24 channels} (raw + \texttt{\_roc} +
\texttt{\_roll\_mean\_30} + \texttt{\_roll\_std\_30}) giving each model explicit
trend, level, and volatility context.

\subsection{Train / Validation Split}

Splits are performed at the \textit{session level} so all windows from a given
session stay together, preventing temporally correlated windows from leaking
across the boundary. An 80/20 stratified split was used across the 16 training
sessions; 2 sessions were held out as unseen test data throughout the entire study.

% =============================================================================
\section{Model Selection}
% =============================================================================

Four model families were evaluated. Each was chosen to test a different
trade-off between interpretability, temporal modelling capacity, and
suitability for embedded deployment.

\begin{table}[H]
\centering
\caption{Model characteristics.}
\begin{tabularx}{\textwidth}{lXX}
\toprule
\textbf{Model} & \textbf{Strengths} & \textbf{Limitations} \\
\midrule
\textbf{Random Forest} &
    Interpretable feature importances; robust on small tabular datasets;
    no scaling required. &
    Requires hand-crafted window aggregates; does not learn from raw
    time-series directly. \\
\addlinespace
\textbf{XGBoost} &
    Gradient boosting corrects residual errors sequentially, typically
    outperforming RF on the same features; built-in regularisation. &
    Same tabular feature dependency as RF; less interpretable than RF. \\
\addlinespace
\textbf{1D CNN} &
    Learns temporal patterns directly from raw signal; no manual feature
    engineering needed; parallelisable and fast to train. &
    Requires more data to generalise well; local receptive field
    may miss longer-range dependencies. \\
\addlinespace
\textbf{TCN} &
    Dilated causal convolutions give an exponentially wider receptive field
    with fewer parameters than stacked Conv1D; residual connections
    stabilise training. &
    Slightly more complex to tune than CNN; still data-hungry compared
    to tree-based methods. \\
\bottomrule
\end{tabularx}
\end{table}

Feature importances from the Random Forest identified HCHO and RH as the
dominant discriminators between classes, which informed the channel selection
for the CNN and TCN input.

% =============================================================================
\section{Results}
% =============================================================================

\subsection{Validation Accuracy}

\begin{table}[H]
\centering
\caption{Validation accuracy on held-in sessions (115 windows).}
\begin{tabular}{llr}
\toprule
\textbf{Run} & \textbf{Model} & \textbf{Val acc} \\
\midrule
Run 1 & RF (50 features)    & 99.1\% \\
Run 2 & RF (top-20)         & 98.3\% \\
Run 3 & RF (all 288)        & 99.1\% \\
Run 7 & XGBoost             & 99.1\% \\
Run 4 & CNN R1              & 99.4\% \\
Run 5 & CNN R2              & 99.6\% \\
Run 6 & TCN R1              & 99.5\% \\
Run 6 & TCN R2              & 99.6\% \\
\bottomrule
\end{tabular}
\end{table}

\subsection{Test Accuracy on Held-Out Sessions}

\begin{table}[H]
\centering
\caption{Test accuracy on \texttt{blood\_9} and \texttt{sweat\_7} (never seen during training).}
\begin{tabular}{llrr}
\toprule
\textbf{Run} & \textbf{Model} & \textbf{blood\_9} & \textbf{sweat\_7} \\
\midrule
Run 1 & RF — 50 features    & 100.0\% & 87.9\% \\
Run 2 & RF — top-20         & 100.0\% & 89.7\% \\
Run 3 & RF — all 288        & 100.0\% & 100.0\% \\
Run 7 & XGBoost             & 100.0\% & 100.0\% \\
Run 4 & CNN R1 (3-layer)    & 99.3\%  & 100.0\% \\
\textbf{Run 5} & \textbf{CNN R2 (2-layer)} & \acc{100.0\%} & \acc{100.0\%} \\
Run 6 & TCN R1              & 98.5\%  & 100.0\% \\
Run 6 & TCN R2 (reduced)    & 99.3\%  & 96.6\% \\
\bottomrule
\end{tabular}
\end{table}

CNN R2 and XGBoost achieve the best overall result: 100\% accuracy on both
held-out test sessions. RF Run 1 underperforms on \texttt{sweat\_7} (87.9\%),
likely because the feature selection threshold discards several sweat-relevant
features; adding all features (Run 3) recovers the gap.

% =============================================================================
\section{Sanity Check}
% =============================================================================

Two tests were conducted to verify that strong test results reflect genuine
learned patterns rather than a lucky data split or structural data artefacts.

\subsection{Test 1 — Seed Sweep}

The session split was repeated with 5 random seeds \{42, 0, 7, 123, 99\}. CNN R2
and TCN R2 were retrained from scratch for each seed and evaluated on the fixed
test sessions.

\begin{table}[H]
\centering
\caption{Seed sweep results — mean $\pm$ std across 5 seeds.}
\begin{tabular}{llrrrr}
\toprule
\textbf{Model} & \textbf{Session} & \textbf{Mean} & \textbf{Std} & \textbf{Min} & \textbf{Max} \\
\midrule
CNN R2 & blood\_9 & 99.4\% & 1.2\% & 97.0\% & 100.0\% \\
CNN R2 & sweat\_7 & 99.7\% & 0.7\% & 98.3\% & 100.0\% \\
TCN R2 & blood\_9 & 99.4\% & 0.6\% & 98.5\% & 100.0\% \\
TCN R2 & sweat\_7 & 99.5\% & 0.7\% & 98.3\% & 100.0\% \\
\bottomrule
\end{tabular}
\end{table}

\textbf{Verdict: \pass.} Standard deviation $\leq 1.2\%$ across all seeds.
Results are not dependent on a particular train/val split.

\subsection{Test 2 — Permutation Test}

Training labels were randomly shuffled 5 times. A new model was fitted on
scrambled labels each time and evaluated on the real test labels.

\begin{table}[H]
\centering
\caption{Permutation test results — permuted mean vs real accuracy.}
\begin{tabular}{llrrrl}
\toprule
\textbf{Model} & \textbf{Session} & \textbf{Perm mean} & \textbf{Real acc} & \textbf{Gap} & \textbf{Verdict} \\
\midrule
CNN R2 & blood\_9 & 34.1\% & 100.0\% & +65.9\% & \pass \\
CNN R2 & sweat\_7 & 49.7\% & 100.0\% & +50.3\% & \pass \\
TCN R2 & blood\_9 & 22.5\% &  99.3\% & +76.7\% & \pass \\
TCN R2 & sweat\_7 & 36.2\% &  96.6\% & +60.3\% & \pass \\
\bottomrule
\end{tabular}
\end{table}

\textbf{Verdict: \pass.} Permuted accuracy collapses to near-random ($\sim$33\%
expected for 3 classes). The real models outperform by 50–77 percentage points,
confirming they learn genuine gas signal patterns rather than exploiting
session structure or data artefacts.

\end{document}
